% context: 9th-grade physics homework, Unit 2 kinematics, Giancoli 2-4 and 2-5
% topic: simple calculations at constant acceleration with the four kinematic
%        equations; speeding up from rest, braking to a stop, a negative acceleration;
%        the Vernier cart from 2-2CW returns (its data fit x = 0.25 t^2)
% structure: two pages with room to work; page 1 the equations, the cart, and a
%        fast car; page 2 braking distance and a skyscraper elevator. No free fall yet.
% data: Tesla's advertised 0-60 mi/h time for the Model S Plaid; braking distances
%        from the UK Highway Code table; Shanghai Tower elevator top speed and height
\documentclass[12pt, twoside]{article}
%\documentclass[12pt, twoside]{article}
\usepackage[letterpaper, margin=1in, headsep=0.2in]{geometry}
\setlength{\headheight}{0.6in}
%\usepackage[english]{babel}
\usepackage[utf8]{inputenc}
\usepackage{microtype}
\usepackage{amsmath}
\usepackage{amssymb}
%\usepackage{amsfonts}
\usepackage[nomessages]{fp} %\FPeval{\var-name}{2*sin(pi/6)}
\usepackage{siunitx} %units in math. eg 20\milli\meter
\usepackage{yhmath} % for arcs, overparenth command
\usepackage{tikz} %graphics
\usetikzlibrary{quotes, angles, arrows, arrows.meta}
\usepackage{graphicx} %consider setting \graphicspath{{images/}}
\usepackage{parskip} %no paragraph indent
\usepackage{enumitem}
\usepackage{multicol}
\usepackage{venndiagram}

\usepackage{fancyhdr}
\pagestyle{fancy}
\fancyhf{}
\renewcommand{\headrulewidth}{0pt} % disable the underline of the header
\raggedbottom
\hfuzz=2mm %suppresses overfull box warnings

\usepackage{hyperref}
\usepackage{float}

\title{Physics}
\author{Chris Huson}
\date{October 2026}

\fancyhead[LE]{\thepage}
\fancyhead[RO]{\thepage \\ Nickname: \hspace{2.25cm} \,\\ \,}
\fancyhead[LO]{La Scuola d'Italia / Huson / Physics \\* 23 October 2026}

\begin{document}

\subsubsection*{2.3 Homework: Constant Acceleration}

\textbf{Reference (provided).} \quad These hold when the acceleration $a$ is constant.

\vspace{0.1cm}
{\hspace{0.5cm}%
\renewcommand{\arraystretch}{1.8}%
\begin{tabular}{l @{\hspace{2cm}} l}
  $v = v_0 + at$ & $v^2 = v_0^2 + 2a(x - x_0)$ \\
  $x = x_0 + v_0 t + \frac{1}{2}at^2$ & $\bar{v} = \dfrac{v + v_0}{2}$
\end{tabular}}

\vspace{0.2cm}
For each problem, list what you know ($v_0$, $v$, $a$, $t$, $x - x_0$) and what you
want, then choose the equation that connects them. Show the units.

\vspace{0.2cm}

\begin{enumerate}[itemsep=0.5cm]

\item The cart from 2.2 is released from rest on a tilted Vernier track. It speeds
  up at a constant $a = 0.500$ m/s$^2$.
  \begin{enumerate}[itemsep=0.4cm]
    \item An acceleration of 0.500 m/s$^2$ means the velocity grows by 0.500 m/s
      every second. Fill in the table.

      \vspace{0.1cm}
      \renewcommand{\arraystretch}{1.5}
      \begin{tabular}{| c | p{1.3cm} | p{1.3cm} | p{1.3cm} | p{1.3cm} | p{1.3cm} |}
        \hline
        $t$ (s) & \centering 0 & \centering 1.00 & \centering 2.00 &
          \centering 3.00 & \centering\arraybackslash 4.00 \\ \hline
        $v$ (m/s) & & & & & \\ \hline
      \end{tabular}
    \item How far has the cart rolled at $t = 2.00$ s? In 2.2, Graphical Analysis
      measured 1.000 m. Does your answer agree?
      \vspace{3.0cm}
  \end{enumerate}

\item Tesla advertises that its Model S Plaid goes from 0 to 60 mi/h in 1.99 s.
  Take 60 mi/h $= 26.8$ m/s and assume the acceleration is constant.
  \begin{enumerate}[itemsep=0.4cm]
    \item Find the car's acceleration.
      \vspace{3.0cm}
    \item How far does the car travel in those 1.99 s?
  \end{enumerate}

\newpage

\item The UK Highway Code, the official rules for British drivers, says that a car
  at 30 mi/h (13.4 m/s) needs 14 m of braking distance to stop.
  \begin{enumerate}[itemsep=0.4cm]
    \item Find the car's acceleration while it brakes. Why is it negative?
      \vspace{3.0cm}
    \item At 60 mi/h (26.8 m/s), with the same acceleration, find the braking
      distance. Twice the speed gives how many times the braking distance?
      \vspace{3.0cm}
  \end{enumerate}

\item The elevators in the Shanghai Tower, 632 m tall, reach a top speed of
  20.5 m/s. Suppose an elevator starts from rest and speeds up at a constant
  1.0 m/s$^2$, about as much as passengers find comfortable.
  \begin{enumerate}[itemsep=0.4cm]
    \item How long does it take to reach top speed?
      \vspace{2.3cm}
    \item How far does it rise while speeding up?
      \vspace{2.8cm}
    \item At the top it slows down at the same rate, so it needs the same distance
      to stop. What is the shortest ride on which the elevator reaches top speed?
  \end{enumerate}

\end{enumerate}
\end{document}
